\documentclass[5pt, a4paper]{article}

\usepackage[utf8]{inputenc}
\usepackage[T1]{fontenc}
\usepackage{textcomp}
\usepackage{amsmath, amssymb}
\usepackage{multicol}
\usepackage{proof}
\usepackage{enumitem}
\usepackage{fancyhdr}
\usepackage[margin=1in]{geometry}
% figure support
\usepackage{import}
\usepackage{xifthen}
\pdfminorversion=7

\usepackage[most]{tcolorbox}
\usepackage{pdfpages}
\usepackage{transparent}
\newcommand{\incfig}[1]{%
    \def\svgwidth{\columnwidth}
    \import{./figures/}{#1.pdf_tex}
}

\tcbset{
colback=gray!5,
colframe=blue!50!white,
coltitle=black,
arc=0pt,
boxrule=-0.1mm,
fonttitle=\bfseries,
coltitle=black,
left=2.5mm,
leftrule=1mm,
right=3.5mm,
colbacktitle=black!10!white,
toptitle=0.75mm,
bottomtitle=0.5mm,
}
\begin{document}
\begin{tcolorbox}[title=Question 1]
    Verify the expressions \[
    \cos \left( \frac{(k + \frac{1}{2})\pi}{n+1} \right) 
    , k = 0,1,2,\ldots,n\] 
    and \[
    \cos \left( \frac{(2k-1)\pi}{2m} \right) 
    , k = 1,2,3,\ldots,m\] 
    determine the same set of values when $m = n + 1$
\end{tcolorbox}
Suppose some $m = n + 1$. Then,
\[
\cos \left( \frac{(2k + 1)\pi}{2m} \right) =
\cos \left( \frac{(2k-1)\pi}{2m} \right) 
.\] 

We'll reindex the LHS. Define $j = k + 1$, such that the expression is
preserved. We'll have
\[
\cos \left( \frac{(j - 1)\pi}{m} \right) =
\cos \left( \frac{(2k-1)\pi}{2m} \right) 
.\] 

The LHS is defined for values of $j = 1,2,3,\ldots,n+1$. The RHS is also defined
for values of $k = 1,2,3,\ldots,m$. Since $m=n+1$, both of these expressions
represents the same set of values.
\begin{tcolorbox}[title=Question 2]
    Use calculus to verify that the first three Legendre polynomials are
    orthogonal to eachother.
\end{tcolorbox}
Recall that the first three Legendre polynomial is given by \[
    P_0(x) = 1, P_1(x) = x, P_2(x) = \frac{1}{2}(3x^2 -1)
.\] 
A polynomial is said to be orthogonal it's inner product is zero. \[
    \int_{-1}^{1} f(x)g(x)dx = 0
.\] 
\subsubsection*{$P_0$ and $P_1$}
\begin{align*}
    \int_{-1}^{1} 1 \cdot x dx = \left[\frac{x^2}{2}\right]_{-1}^{1} = 0
.\end{align*}
\subsubsection*{$P_0$ and $P_2$}
\begin{align*}
    \int_{-1}^{1} 1 \cdot \frac{1}{2}(3x^2 - 1) dx &= \frac{1}{2}\left( 3
    \int_{-1}^{1} x^2dx - \int_{-1}^{1} dx  \right) \\
                                                   &= \frac{1}{2}\left[ 3 \left(
                                                   \frac{x^3}{3}\right) - x
                                               \right]_{-1}^{1} \\
                                               &= 0 
.\end{align*}
\subsubsection*{$P_1$ and $P_2$}
\begin{align*}
    \int_{-1}^1 x \cdot \frac{1}{2}(3x^2 - 1) dx 
    &= \frac{1}{2} \int_{-1}^1 (3x^3 - x) dx\\
    &= \frac{1}{2} \left[ 3\int_{-1}^{1}x^3dx - \int_{-1}^{1} xdx  \right] \\
    &= \frac{1}{2} \left[ 3 \left( \frac{x^4}{4} \right) -
    \left(\frac{x^2}{2}\right) \right]_{-1}^{1} = 0
.\end{align*}
\begin{tcolorbox}[title=Question 3]
    The Chebyshev nodes on a general interval $[a,b]$ are given  \[
        x_k = \frac{1+b}{b} + \frac{b - a}{a} \cos \left(
        \frac{(k+\frac{1}{2})\pi}{n+1} \right) 
    ,k = 0,1,2,\ldots,n.\] 
    List the Chebyshev interpolation nodes $x_0,x_1,\ldots,x_n$ in the following
    interval:
    \begin{enumerate}[label=\alph*)]
        \item $[-1,1], n = 5$
        \item $[-2,2], n = 3$
        \item $[2,12], n = 5$
    \end{enumerate}
\end{tcolorbox}
We follow the formula and list the nodes for each interval. For $[-1,1]$, we
have  \[
x_k = \cos\left( \frac{(k+\frac{1}{2})\pi}{6} \right) = \cos\left( \frac{(2k +
1)\pi}{12} \right)
.\] 
Thus, \[
    x_0 = \cos \left( \frac{\pi}{12} \right) 
x_1 = \cos\left( \frac{\pi}{4} \right) x_2 = \cos\left( \frac{5}{12}\pi \right)  
x_3 = \cos\left( \frac{7}{12} \pi  \right) 
x_4 = \cos \left( \frac{9}{12} \pi \right) 
x_5 = \cos \left( \frac{11}{12} \pi \right) 
.\] 

FOr $[-2,2]$ \[
x_k  = 2\cos \left( \frac{(2k+1)\pi}{8} \right) 
.\] 
Thus, 
\[
x_0 = 2\cos \left( \frac{\pi}{8} \right) 
x_1 = 2\cos \left( \frac{3\pi}{8} \right) 
x_2 = 2\cos \left( \frac{5\pi}{8} \right) 
x_3 = 2\cos \left( \frac{7\pi}{8} \right) 
.\] 
For $[2,12]$,  \[
    x_k = 7 + 5 \cos \left( \frac{(2k + 1) \pi}{12}\right) 
.\] 
Thus, \[
x_0 = 7 + 5 \cos\left( \frac{\pi}{12} \right) 
x_1 = 7 + 5 \cos\left( \frac{3\pi}{12} \right) 
x_2 = 7 + 5 \cos\left( \frac{5\pi}{12} \right) 
\]
\[
x_3 = 7 + 5 \cos\left( \frac{7\pi}{12} \right) 
x_4 = 7 + 5 \cos\left( \frac{9\pi}{12} \right) 
x_5 = 7 + 5 \cos\left( \frac{11\pi}{12} \right) 
.\] 

\begin{tcolorbox}[title=Question 4]
    Consider the data points $(-1,1),(1,2),(2,3).$ 
    \begin{enumerate}[label=\alph*)]
        \item Write down the system of linear equations that determines the
            piecewise quadratic interpolant that is twice differentiable at $x =
            1$
        \item Wirte down the linear equation that determines the natural cubic
            splines interoplant for these data.
    \end{enumerate}
\end{tcolorbox}
\begin{tcolorbox}[title=Question 5]
    Let $T_0,\ldots,T_n$ denote the Chebyshev polynomials.
    \begin{enumerate}[label=\alph*)]
        \item Compute $T_4,T_5,T_6$ using the recurrence relation.
        \item Verivy that the roots of $T_4$ are claimed.
        \item Using calculus, compute the local extrema of $T_4$
    \end{enumerate}
\end{tcolorbox}
\begin{tcolorbox}[title=Question 6]
    Write down the interpolating polynomial that interpolates $f(x) = \cos x$ on
    the interval $[0,2\pi]$ using 6 Chebyshev nodes. SHow atleast 4 significand
    figures for each coeffecient. 
\end{tcolorbox}
\end{document}
